
recovers the classical ``geometric'' notion of algebraic equivalence considered in 2.4.1. We therefore prefer to say that an element of @@M0@@ is algebraically equivalent to zero when its image in @@M1@@ is so, where @@M2@@ is an algebraic closure of @@M3@@, and to proceed in the same way for the definition of algebraic equivalence in the groups @@M4@@, @@M5@@ and @@M6@@.

\sgaStableTarget{sga6.e10.para.7-12}\textbf{7.12.} With these definitions in place, I claim that \emph{the specialization homomorphism} (7.3.1) \emph{carries elements algebraically equivalent to zero to elements algebraically equivalent to zero}, and consequently defines, by passage to the quotient, a homomorphism
\sgaStableTarget{sga6.e10.eq.7-12-1}@@M7@@

To see this, one reduces immediately to the case where @@M8@@ is complete and @@M9@@ is algebraically closed (using a straightforward compatibility made explicit in 7.15). It suffices to see that @@M10@@ carries an element of the form (7.11.1) of @@M11@@ to an element algebraically equivalent to zero in @@M12@@. Here @@M13@@, and we shall accordingly write @@M14@@, @@M15@@ in place of @@M16@@, @@M17@@ above. One knows that @@M18@@ is projective over @@M19@@, so (EGA IV 2.8.5) there is a projective flat scheme @@M20@@ over @@M21@@ whose generic fiber is @@M22@@. Using Abhyankar's resolution of singularities of arithmetic surfaces [6], one sees that one can even choose @@M23@@ to be a regular scheme. (N.B. Flatness is preserved under resolution, by EGA IV 14.3.8.) Since @@M24@@ is proper over @@M25@@, the points @@M26@@ and @@M27@@ of @@M28@@ rational over @@M29@@ extend to sections @@M30@@, @@M31@@ of @@M32@@ over @@M33@@, and, since @@M34@@ is regular, it is well known that @@M35@@ is smooth over @@M36@@ at the points of @@M37@@, @@M38@@ (EGA IV, 17.12.1 c) @@M39@@ a) and 19.1.1). As a 